\documentclass[11pt]{article}

\usepackage[margin=1in]{geometry}
\usepackage{microtype}
\usepackage{amsmath}
\usepackage{graphicx}
\usepackage{booktabs}
\usepackage[hidelinks]{hyperref}

\graphicspath{{figures/}}

\newcommand{\natex}{\texttt{natex}}

\title{A Break Made of Growth: The 2025Q1 Datacenter Surge as Smooth
Exponential Expansion, Not a Stargate Regime Change}
\author{Hauke Hillebrandt\thanks{University College London. Email:
\texttt{ucjthhi@ucl.ac.uk}. This paper and the underlying \natex{} software were
prepared with substantial assistance from Anthropic's Claude models; the author
reviewed the analyses and text and is responsible for all remaining errors.}}
\date{July 2026 \\ \small analysis pass of record: 2026-07-19, seeded
(\texttt{--seed 0}), \natex{} v0.2.0 at commit \texttt{cce6186}}

\begin{document}

\maketitle

\begin{abstract}
\noindent In Epoch AI's Frontier Data Centers panel (73 sites, 2023Q4--2026Q2),
tracked quarterly IT-power additions rise roughly fivefold after 2025Q1 --- the
quarter of the Stargate announcement and DeepSeek-R1 --- from a 296.8~MW/q
pre-period mean to 1{,}516.9~MW/q, with complete pre/post separation (minimum
post quarter 537~MW $>$ maximum pre quarter 448~MW). We show this
``Stargate-quarter industry break'' is a correctly-identified artifact, and
that attribution to ``neocloud'' operators is a clean null. On a placebo
break-date grid (all 6 feasible dates), 2025Q1 is \emph{never} the best break
date --- rank 3--6 of 6 on all three statistics, with 2025Q4 the argmax on two
--- and log tracked capacity grows near-linearly with mildly
\emph{decelerating} growth ($0.40 \to 0.33$ log/q); the log-level kink is
negative at every candidate date. The levels ``break'' is what smooth
$\sim$40\%-per-quarter exponential growth looks like. A seeded \natex{} SuDDDS
scan does return $t_0 = 2025$Q1 \emph{exactly} with scan $p=0.010$, but
discovery reads only the design columns $(x, t, \theta)$, never the outcome:
the hit mechanically recovers the constructed 21-of-73-site treatment coding
and contributes no outcome evidence --- a role-assignment triage that corrects
the scout pass's interpretation. Attribution fails directly: a two-way
fixed-effects contrast of neocloud versus other sites gives
$\tau = +0.098$ asinh-MW ($t=0.699$) with a 199-draw label-permutation
$p=0.56$. The surge is instead broad-based and hyperscaler-led: the neocloud
group carries 23.7\% of post-2025Q1 additions, and the top site is the
Anthropic--Amazon New Carlisle campus at 10.0\%. All conclusions survive
panel-end truncation and an incumbent-only variant, and the panel is provably
invariant to the projection-snapshot filter.
\end{abstract}

\section{Introduction}\label{sec:intro}

January 2025 was the most crowded month in the short history of AI
infrastructure policy: the US diffusion rule was published on 13~January,
DeepSeek-R1 was released on 20~January \cite{deepseekr1}, and the \$500-billion
Stargate project was announced on 21~January \cite{openai2025}. A natural
reading of subsequent buildout data is that 2025Q1 marks a regime change ---
the quarter the industry shifted from ordinary expansion to a land rush --- and
that the shift is concentrated in the ``neocloud'' operators (CoreWeave,
xAI/Colossus, the Stargate sites, and peers) whose business is AI compute
rather than general cloud. This paper tests both claims at the site level and
rejects both, in an instructive way: the industry-wide ``break'' is real in
levels but is not a break, and the neocloud attribution is a clean null.

The buildout itself is well documented. Data centers are the engine rooms of
AI: industrial facilities whose construction, power, and cooling determine
where frontier compute can exist \cite{pilzheim2023}. The leading AI
supercomputers have doubled in performance roughly every nine months, with
hardware cost and power needs doubling yearly \cite{pilz2025}, and the
International Energy Agency projects global data-centre electricity demand to
more than double by 2030, driven chiefly by AI \cite{iea2025}. What has been
missing is site-level measurement: Epoch AI's Frontier Data Centers hub
\cite{epochdc} now tracks individual campuses from satellite imagery, permits,
and disclosures, making the timing question empirically addressable.

Methodologically, the exercise sits where two cautionary literatures meet.
Breaks located in \emph{time} lack the cross-sectional randomization logic of
standard quasi-experimental designs, and naive inference on a single temporal
cutoff is fragile \cite{hausmanrapson2018}; when identification comes from one
calendar event hitting one group, asymptotic cluster inference is unreliable
and permutation-style inference over placebo assignments is the appropriate
discipline \cite{conleytaber2011}. We add a third, automation-specific
caution: when a discovery algorithm is handed a panel whose treatment coding
was \emph{constructed} from the hypothesis under test, an exact recovery of
the hypothesized date is confirmation of the coding, not of the outcome. The
scout pass that preceded this study made exactly that misreading; the
role-assignment triage below corrects it.

\section{Data}\label{sec:data}

The input is Epoch AI's Frontier Data Centers hub timelines \cite{epochdc}:
411 dated site snapshots covering 73 sites (2018-12-18 to 2030-01-01),
extracted 2026-07-19. The 73 rows dated on or after 2026-07-25 are stated
projections and are excluded; Section~\ref{sec:results} shows this filter is
provably a no-op for the analysis panel. Each site's ``IT power (MW)'' is
carried forward (LOCF) as a step function evaluated at quarter ends 2023Q4 to
2026Q2, with a site at 0~MW before its first snapshot; quarterly
\emph{additions} are first differences with negative revisions floored at 0;
the site-level outcome is $\operatorname{asinh}(\text{additions}/10)$. Time is
indexed $t = 4\,\text{year} + (\text{quarter}-1)$, so $2025\text{Q1} = 8100$.
The neocloud group comprises 21 sites with hypothesized treatment onset
2025Q1: Colossus~1/2, seven CoreWeave sites, Fluidstack Lake Mariner, eight
OpenAI Stargate sites, Oracle Batam, Start Campus Sines, and VNET Bayin
Ulanqab. The known-treatment comparison panel codes Colossus~1 alone as
treated. Baseline panels are bitwise identical to the scout pass's.

\section{Design and Methods}\label{sec:methods}

\paragraph{Seeded discovery scans.} Two \natex{} \cite{natex} SuDDDS
difference-in-differences scans (Bernoulli model, within-cluster-copula
inference, $q=99$ permutations, seed 0) search the panels for the
best-supported treatment date $t_0$ and window $W$: one on the neocloud
coding, one on the Colossus known-treatment coding. Crucially, SuDDDS
discovery reads only the design columns --- unit, time, and the treatment
indicator $\theta$ --- and never the outcome. Both panels \emph{construct}
$\theta = \text{group} \times \mathbf{1}[t \geq t_0]$, so an exact recovery of
$t_0$ certifies that the constructed assignment is scannable, not that MW
additions break there. Outcome evidence must come from the two designs below.

\paragraph{Aggregate break-date placebos.} On total tracked additions we
compute, at every feasible break date $c$ (six dates with at least three
quarters per side), three statistics: S1, the post$-$pre shift in mean log
additions; S2, the slope change (kink) in the MW-level series at $c$; and S3,
the kink in \emph{log} capacity at $c$, i.e.\ growth acceleration. S2 and S3
are single-series analogues of the kink contrasts formalized in the
difference-in-kinks literature \cite{bockerman2025} and implemented in
\natex{} \cite{natex}. If 2025Q1 is a genuine regime break it should be the
argmax of the grid; its rank across placebo dates is the test, in the spirit
of randomization inference over assignments \cite{conleytaber2011}, sized for
a coarse grid (placebo-$p$ floor $1/6$).

\paragraph{Neocloud attribution.} A two-way fixed-effects (TWFE) regression of
$\operatorname{asinh}(\text{additions}/10)$ on
$\text{neocloud} \times \mathbf{1}[t \geq 2025\text{Q1}]$ with site and
quarter fixed effects (73 sites $\times$ 11 quarters) estimates the
neocloud-specific post-2025Q1 shift $\tau$. Because the contrast is a single
calendar event hitting one group, inference is by label permutation
\cite{conleytaber2011}: 199 draws (seed 0) reassigning the 21-site label among
all 73 sites, plus a chip-protocol variant reassigning among the 52 controls.
CR1 cluster-by-site standard errors are reported alongside and agree.

\paragraph{Decomposition and sensitivity.} Post-2025Q1 additions are
decomposed by site and group. Sensitivity variants truncate the panel end at
2026Q1 and 2025Q4, and an incumbent-only variant keeps the 59 sites first
tracked on or before 2024-12-31 (12 treated) to separate real expansion from
satellite-coverage entry \cite{epochdc,hausmanrapson2018}.

\section{Results}\label{sec:results}

\begin{figure}[t]
\centering
\includegraphics[width=\textwidth]{fig1}
\caption{(a) Tracked quarterly IT-power additions, stacked by group: the
post-2025Q1 rise is roughly fivefold and occurs in \emph{both} groups; the
largest single quarters are 2025Q4 and 2026Q2, not 2025Q1. (b) Log total
tracked capacity is near-linear ($\sim$0.4 log points per quarter) with mildly
decelerating growth --- the levels ``break'' is smooth exponential growth
read in levels. (c) Placebo break-date grid for S1, the post$-$pre shift in
mean log additions: 2025Q1 (vermilion) ranks 3 of 6; the argmax is 2025Q4.
(d) Top-10 site contributions to post-2025Q1 additions: broad-based and
hyperscaler-led, with the neocloud group (orange) carrying 23.7\% overall.}
\label{fig:fig1}
\end{figure}

\paragraph{The surge is real in levels.} Tracked additions average 296.8~MW/q
over 2023Q4--2024Q4 and 1{,}516.9~MW/q over 2025Q1--2026Q2 --- a ratio of
5.11 --- with complete separation: the smallest post-2025Q1 quarter (537~MW)
exceeds the largest pre-quarter (448~MW). Nothing below disputes this; the
question is whether it is a \emph{localized break at 2025Q1}.

\paragraph{But 2025Q1 is never the best break date.} On the six-date placebo
grid (Table~\ref{tab:main}, Panel~A), 2025Q1 is the argmax of none of the
three statistics. The S1 log-additions shift at 2025Q1 is 1.488, rank 3 of 6
(placebo $p=0.50$); the argmax is 2025Q4 at 1.665. The S2 level kink at
2025Q1 is 1{,}354~MW/q, dead last --- rank 6 of 6 ($p=1.00$), argmax 2025Q4
at 2{,}119. The S3 log-level kink is $-0.100$ at 2025Q1 and \emph{negative at
all six candidate dates}: quarterly log growth of tracked capacity
decelerates from 0.404 to 0.327 log/q. A series growing smoothly at
$\sim$40\% per quarter mechanically makes every later quarter's additions
exceed every earlier one's; the ``break'' is the exponential, not a regime
change. The largest single jumps are 2025Q4 ($+1{,}936$~MW) and 2026Q2
($+3{,}610$~MW). Across all sensitivity variants, 2025Q1 is the argmax in 1
of 12 variant-statistic cells (the 2025Q4-truncated S3, placebo $p=0.25$ on a
four-date grid).

\paragraph{The SuDDDS 2025Q1 hit is mechanical.} The seeded neocloud scan
returns max LLR $=24.114$, $t_0 = 8100 = 2025$Q1 \emph{exactly}, $W=3$, scan
$p=0.010$ (the $q=99$ permutation floor), and passes composition validation.
This is the scout pass's headline --- and it is a recovery of the constructed
treatment coding, not outcome evidence: discovery never reads the outcome
(Section~\ref{sec:methods}), and the scan's in-scan effect estimators refused
in both runs (0 usable placebo units of the required 5), so the scans
contribute no outcome estimate. The Colossus known-treatment scan illustrates
the power ceiling with 1 treated site of 73: LLR $=2.537$, $t_0=2024$Q3,
$W=5$, scan $p=0.230$, deterministic on re-run (the scout logged $p=0.28$
from a marginally earlier code state with identical LLR/$t_0$/$W$).

\paragraph{Neocloud attribution is a clean null.} The TWFE contrast gives
$\tau = +0.0976$ asinh-MW (CR1 SE $=0.1396$, $t=0.699$). The 199-draw label
permutation puts the observed $\tau$ well inside the null: $p=0.560$
relabeling among all 73 sites, $p=0.555$ under the chip protocol, with a
permutation $\tau$ SD of 0.150 --- the estimate is $\sim$0.65 null SDs from
zero. Truncating the panel at 2026Q1 or 2025Q4 leaves the null intact
($p=0.605$, $p=0.505$). The incumbent-only variant is larger but still not
significant: $\tau=+0.258$, $t=1.49$, permutation $p=0.155$.

\paragraph{Broad-based and hyperscaler-led.} Post-2025Q1 additions total
9{,}101~MW, of which the 21-site neocloud group carries 2{,}158~MW ---
23.7\%. Concentration is moderate: the top 5 sites carry 33.9\% and the top
10 carry 49.1\%. The leaders are the Anthropic--Amazon New Carlisle campus at
910~MW (10.0\%), Microsoft Fairwater Atlanta (636), Meta Prometheus (631),
Colossus~2 (490, neocloud), and OpenAI Stargate Abilene (421, neocloud).
Inside the scan window 2025Q1--2025Q3 (2{,}245~MW) the largest contributors
are New Carlisle (398~MW, 17.7\%) and Amazon Madison (284~MW, 12.7\%) ---
hyperscaler campuses, not neoclouds.

\paragraph{Sensitivity and the projection filter.} Excluding stated
projections is provably a no-op: all 73 projection rows are dated on or after
2026-07-25, later than the last quarter-end 2026-06-30, and every site has a
real snapshot on or before 2026-06-30, so the 2023Q4--2026Q2 panel is bitwise
invariant to any snapshot-filter date from 2026-06-30 on (including no
filter). The incumbent variant shows the surge is not satellite-coverage
entry: the 59 sites first tracked by 2024-12-31 carry 3{,}074 of the
3{,}610~MW of 2026Q2 additions ($\sim$85\%).

\begin{table}[t]
\centering
\footnotesize
\setlength{\tabcolsep}{3pt}
\caption{Headline estimates. Panel A: aggregate break-date placebo grid (six
feasible dates, $\geq 3$ quarters per side); the test is 2025Q1's rank.
Panel B: seeded \natex{} SuDDDS scans (Bernoulli/WCC, $q=99$, seed 0).
Panel C: TWFE neocloud attribution with 199-draw label-permutation inference
(seed 0).}
\label{tab:main}
\begin{tabular}{lrrrl}
\toprule
Quantity & Estimate & Rank / $t$ & $p$ & notes \\
\midrule
\multicolumn{5}{l}{\emph{Panel A: is 2025Q1 the best break date? (placebo grid, $n=6$ dates)}}\\
S1 shift in mean log additions & $1.488$ & 3/6 & $0.50$ & argmax 2025Q4 ($1.665$) \\
S2 level kink (MW/q)           & $1{,}354$ & 6/6 & $1.00$ & argmax 2025Q4 ($2{,}119$) \\
S3 log-level kink (log/q)      & $-0.100$ & 4/6 & $0.67$ & negative at all 6 dates \\
post/pre mean additions ratio  & $5.11$ & --- & --- & min post $537 >$ max pre $448$ \\
log growth, pre $\to$ post     & $0.404 \to 0.327$ & --- & --- & deceleration \\
\midrule
\multicolumn{5}{l}{\emph{Panel B: SuDDDS scans (discovery reads $(x,t,\theta)$ only --- never the outcome)}}\\
neocloud scan, max LLR    & $24.114$ & $t_0=2025$Q1, $W=3$ & $0.010$ & recovers the coding \\
Colossus scan, max LLR    & $2.537$  & $t_0=2024$Q3, $W=5$ & $0.230$ & 1 treated of 73 \\
\midrule
\multicolumn{5}{l}{\emph{Panel C: neocloud attribution, TWFE $\tau$ (asinh-MW), site+quarter FE}}\\
baseline (21/73 treated)  & $+0.098$ & $t=0.699$ & $0.560$ & $\tau$ SD $0.150$; chip $p=0.555$ \\
truncate 2026Q1 / 2025Q4  & $0.088/0.112$ & --- & $0.605/0.505$ & conclusions unchanged \\
incumbent-only (12/59 treated) & $+0.258$ & $t=1.49$ & $0.155$ & larger, still n.s. \\
\bottomrule
\end{tabular}

\medskip
\parbox{0.92\textwidth}{\footnotesize Notes: Panel A placebo $p$ is the rank
of 2025Q1 divided by the grid size (floor $1/6$). Panel B scan $p$-values are
$q=99$ within-cluster-copula permutation floors; the neocloud scan's exact
$t_0$ recovery certifies the constructed 21/73-site assignment is scannable,
not that the outcome breaks there (in-scan effect estimators refused: 0
usable placebos of $\geq 5$ required). Panel C $p$-values are 199-draw label
permutations, seed 0; CR1 cluster-by-site $t$-statistics agree with the
permutation verdicts.}
\end{table}

\section{Caveats and Conclusion}\label{sec:conclusion}

Six caveats bound the claims. First, \emph{LOCF timing}: snapshots are sparse
(median 4 per site), so carried-forward power smears true addition dates by
up to one inter-snapshot gap; quarter-level timing is approximate, and
non-monotone estimate revisions are floored at 0. Second,
\emph{tracked-universe growth}: the hub's coverage grows from 4 to 54 sites
with nonzero power over the panel, so aggregate levels lower-bound industry
capacity and early growth rates are inflated by coverage catch-up --- a bias
\emph{toward} finding late acceleration, which strengthens the
no-2025Q1-break reading, though absolute MW shares remain conditional on
Epoch's tracking choices \cite{epochdc}. Third, \emph{coarse placebo grids}:
with 4--6 feasible dates the placebo $p$ has a floor of $1/6$--$1/4$;
statements are about rank, not tight $p$-values. Fourth, a \emph{crowded
quarter}: Stargate (21~January), DeepSeek-R1 (20~January), and the diffusion
rule (13~January) all land in 2025Q1 \cite{openai2025,deepseekr1}, so even a
real 2025Q1 break could not separate these candidate causes. Fifth,
\emph{calendar-time inference}: the neocloud contrast is one calendar event
hitting one group; we therefore rest on permutation-in-space inference
\cite{conleytaber2011}, and the asymptotic CR1 statistics happen to agree.
Sixth, \emph{scale}: $\operatorname{asinh}(\text{additions}/10)$ compresses
large sites, and $\tau$ lives on that scale.

The conclusion is a double negative with positive content. The post-2025Q1
surge in tracked datacenter additions is real, large, and completely
separated in levels --- but it is what smooth $\sim$40\%-per-quarter
exponential growth of tracked capacity looks like, with 2025Q1 never the best
break date on any statistic, growth mildly decelerating, and the largest
jumps arriving in 2025Q4 and 2026Q2. The one seemingly sharp result --- a
seeded discovery scan returning 2025Q1 exactly at $p=0.010$ --- is a
mechanical recovery of the constructed treatment coding, a failure mode worth
naming for any pipeline that scans panels whose design columns encode the
hypothesis. And the group story fails directly: neocloud attribution is a
clean permutation null, with the buildout led by hyperscaler campuses ---
first among them the Anthropic--Amazon New Carlisle site at 10\% of all
post-2025Q1 additions. The Stargate quarter did not change the growth regime;
the regime was already exponential \cite{pilz2025,iea2025}.

\paragraph{Reproducibility.} All estimates derive from the public Epoch AI
Frontier Data Centers hub timelines \cite{epochdc}; no dataset files are
committed. The only stochastic steps are the seeded scan and permutation
draws (seed 0 throughout). Figure~\ref{fig:fig1} regenerates from the
committed \texttt{figures/make\_fig.py}, which asserts the headline estimates
against the numbers of record before drawing.

\footnotesize
\begin{thebibliography}{9}
\setlength{\itemsep}{1pt}\setlength{\parskip}{0pt}

\bibitem{bockerman2025}
B\"ockerman, P., Jysm\"a, S., and Kanninen, O. (2025).
\emph{Difference-in-Kinks Design}.
IZA Discussion Paper No.~18313.
\url{https://docs.iza.org/dp18313.pdf}

\bibitem{conleytaber2011}
Conley, T.~G., and Taber, C.~R. (2011).
Inference with ``difference in differences'' with a small number of policy
changes.
\emph{Review of Economics and Statistics}, 93(1), 113--125.

\bibitem{deepseekr1}
DeepSeek-AI (2025).
DeepSeek-R1: Incentivizing reasoning capability in LLMs via reinforcement
learning.
arXiv:2501.12948.

\bibitem{epochdc}
Epoch AI (2026).
Frontier Data Centers Hub.
\url{https://epoch.ai/data/data-centers} (timelines extracted 19 July 2026).

\bibitem{hausmanrapson2018}
Hausman, C., and Rapson, D.~S. (2018).
Regression discontinuity in time: Considerations for empirical applications.
\emph{Annual Review of Resource Economics}, 10, 533--552.

\bibitem{natex}
Hillebrandt, H. (2026).
\emph{natex: automated natural-experiment discovery and estimation}.
Software, version 0.2.0.
\url{https://github.com/HaukeHillebrandt/natex}

\bibitem{iea2025}
International Energy Agency (2025).
\emph{Energy and AI}.
IEA, Paris. \url{https://www.iea.org/reports/energy-and-ai}

\bibitem{openai2025}
OpenAI (2025).
Announcing The Stargate Project.
Company announcement, 21 January 2025.
\url{https://openai.com/index/announcing-the-stargate-project/}

\bibitem{pilzheim2023}
Pilz, K., and Heim, L. (2023).
Compute at scale: A broad investigation into the data center industry.
arXiv:2311.02651.

\bibitem{pilz2025}
Pilz, K.~F., Sanders, J., Rahman, R., and Heim, L. (2025).
Trends in AI supercomputers.
arXiv:2504.16026.

\end{thebibliography}

\end{document}
